\section{Conclusion}\label{sec:conclusion}

Nearly all of the results of this paper follow from the spin of gravity and
the two small numbers of Section~\ref{sec:regimes}. Gravity has spin two, so a reading at a single
point holds ten numbers: four principal null directions, a rest frame, and
room for four waves. Matter is slow, so
everything near us is near field: seen at once, where it is now, through any
wall, and never sharply. Matter is tenuous, so nothing bends, reflects or stops
a wave: warm matter glows all the way through, too faintly for anything made
of matter to detect, and every image is a computation.

Each channel has scenes that it cannot tell apart, and they can be stated as
exactly as what it sees: the harmonic part of the density and
nothing more, the swirl of a flow and not its rolls, the shear of a stress and
not its pressure, friction times a power of the temperature and not either
alone, sixty-eight numbers of the sky and not a sixty-ninth. The ambiguities
are removed across channels: by spin, by heat, by an applied tide, by a second
pixel. The most important of them, which hides the interior of a body from its
exterior field, holds for a reading at one time only. A body that is warm, or
to which a tide can be applied, reveals in time
where all its mass is; a black hole does not (Section~\ref{sec:hair}).

The field of light at a point is
a vector and that of gravity a square. From that
difference a pixel reads its own velocity through an unknown room, detects
whether anything in the room is moving, and separates four waves
where a sensor of light separates two. The blur of a near-field image does not
limit the precision of a measurement. A near-field picture holds few bits, tens to hundreds
for each patch of the scene as wide as it is far, and a scene that is simple,
or merely positive, spends them all on a few numbers: the readings that could
never draw a face place a mass to within a millimetre.

The camera of this paper cannot be built, but parts of it exist.
Gradiometers already map what is buried \cite{Stray2022}, a satellite has
flown one around the Earth \cite{RummelYiStummer2011}, and every gravitational-wave observatory models the
pull of its surroundings in order to remove it \cite{Harms2019}. That pull is
the near-field signal described here.
